\section{Problem Description}

We consider a weekly surgery scheduling problem in a hospital operating theatre 
consisting of multiple ORs. At the start of each week, all 
inpatients $p \in P$ on the waiting list must be scheduled for surgery. The 
scheduling process involves three interrelated decisions: (i)~\textit{assignment} 
of patients to OR sessions, (ii)~\textit{sequencing} of patients within each OR, 
and (iii)~\textit{timing}, i.e., determining the planned start time $S_p$ 
communicated to the wards before the shift begins.

Each OR opens at $t_{\text{open}} = 8{:}00$ h and closes at $t_{\text{close}}$. 
The first patient in every OR begins surgery at opening time. Subsequent patients 
may not start before their communicated planned start time, as ward nurses are 
assumed to be perfectly punctual. A deterministic changeover time of 10 minutes 
is required between any two consecutive surgeries.

The key source of uncertainty is the surgery duration $d_p$, which is unknown at 
the time of scheduling. Because surgeries cannot start 
before their planned start time but may finish later than expected, delays 
propagate through the sequence within an OR session.

The objective is to minimise the weighted sum of three cost components: the total patient waiting time, the total OR idle time , and the 
total overtime. 

Because $d_p$ is continuous and its realisation is only observed after scheduling, 
an exact deterministic formulation is insufficient. We therefore adopt a 
two-stage stochastic programming framework with SAA~\cite{Kleywegt}: First-stage decisions (assignment, 
sequencing, timing) are made before uncertainty is revealed, while second-stage 
costs (start time, waiting, idle, overtime) are evaluated over a finite scenario set. The goal is to find a schedule that minimises 
expected costs across all scenarios. 